\newpage
\section{Ponto neutro e margem estática}

% Item 8 do roteiro. Valores preliminares abaixo obtidos com o comando st
% em mn 0.85 e a c 0.5053 (sem trimagem). CONFIRMAR nas rodadas finais.

O ponto neutro foi obtido do comando \texttt{st} do AVL na condição de
cruzeiro, resultando em $x_{np} = 30{,}065$ m, ou 73,2\,\%MAC. As margens
estáticas correspondentes às duas posições de CG são
\begin{equation*}
    SM_{\mathrm{fwd}} = \frac{x_{np} - x_{cg,\mathrm{fwd}}}{c_{\mathrm{ref}}}
    = 57{,}8\,\%\mathrm{MAC},
    \qquad
    SM_{\mathrm{aft}} = \frac{x_{np} - x_{cg,\mathrm{aft}}}{c_{\mathrm{ref}}}
    = 34{,}0\,\%\mathrm{MAC}.
\end{equation*}

Os valores são maiores que os estimados pelo \texttt{designTool} no Lab 02,
que situa o ponto neutro em 45,4\,\%MAC (margens de 30,0\% e 6,2\%). A
diferença é esperada e tem origem na fuselagem, que não é modelada no AVL. O
momento desestabilizante do corpo desloca o ponto neutro real para a frente,
e o \texttt{designTool} contabiliza esse efeito por correlação empírica.

% Discutir qual estimativa e mais conservadora para o projeto e registrar
% os numeros finais apos as rodadas definitivas.
